\newpage\section{Conclusion}

In conclusion, the numerical solution of the shallow water equations was implemented using both an unstaggered grid and a staggered grid using the leapfrog scheme. The result showed that staggering the grid gave a much better agreement with the theoretical solution, compared to the unstaggered grid. It was shown that a twice as large grid spacing in the staggering gave the same accuracy as the unstaggered grid with normal grid spacing. Thus staggering the grid could be used to save computational space and time, while still maintaining the same accuracy or to improve the accuracy of the numerical scheme. The Robert-Asselin filter showed to lower the error of both the unstaggered and staggered grid schemes,  although the affect on the unstaggered grid scheme was not as significant as staggering the grid. Thus, the role of using staggered grids in numerical schemes is an important tool to improve the accuracy of the model. This explains why staggered grids are more commonly used in weather models. 